\clearpage
\begin{explanationbox}
\textbf{从普通切点到顶点：同一个公式覆盖两种方向}

首页取 $a^2=4,b^2=3$，将 $P(4,3)$ 代入即得 $x-y=1$。
图中橙线在 $P$ 处与蓝色曲线有相同的切线方向。半代换是记忆方式，
但只知道直线经过 $P$ 还不能说明它是切线。
下一节会按	extbf{平方差与割线逼近 $	o$ 导数速证 $	o$ 联立判别式}的顺序证明。

\begin{center}
\textbf{顶点处切线竖直，切点式仍然适用}\par\medskip
\begin{tikzpicture}[x=1.1cm,y=1.1cm,>=Stealth,font=\small]
\draw[->,gray] (-4.9,0)--(4.9,0) node[right] {$x$};
\draw[->,gray] (0,-3.2)--(0,3.4) node[above] {$y$};
\draw[blue,very thick,domain=-1.3:1.3,samples=90] plot ({2*cosh(\x)},{sqrt(3)*sinh(\x)});
\draw[blue,very thick,domain=-1.3:1.3,samples=90] plot ({-2*cosh(\x)},{sqrt(3)*sinh(\x)});
\draw[orange!85!black,thick] (2,-3)--(2,3);
\draw[orange!85!black,thick] (-2,-3)--(-2,3);
\fill (2,0) circle (1.8pt) node[below right=4pt] {$(a,0)$};
\fill (-2,0) circle (1.8pt) node[below left=4pt] {$(-a,0)$};
\node[orange!85!black,above] at (2,3.05) {$x=a$};
\node[orange!85!black,above] at (-2,3.05) {$x=-a$};
\node[below left] at (0,0) {$O$};
\end{tikzpicture}

\end{center}
图取 $a=2,b=\sqrt3$。观察左右顶点：曲线在顶点处的切线都平行于 $y$ 轴。
把 $P=(a,0)$ 或 $P=(-a,0)$ 代入切点式，分别得到 $x=a$ 或 $x=-a$。
因此，\textbf{切点式不需要斜率存在}。

\textbf{理解边界：}双曲线与一条直线只有一个公共点，并不足以保证相切；
下一节通过割线的方向、导数与二次方程的重根分别核对相切，
不能把交点个数当作相切的完整判据。

	extbf{先理解证明主线：}设曲线上有两点 $P(x_0,y_0)$、$Q(x,y)$，
将两点代入原方程后	extbf{相减}，就能用平方差公式得到割线 $PQ$ 的斜率；
当 $Q$ 沿曲线越来越靠近 $P$ 时，割线方向越来越接近切线方向。
这也解释了为什么半代换公式不只是一个记忆口诀。
\end{explanationbox}
